%% =====================================================================
%%  2 Problem Restatement and Clarification
%% =====================================================================
%%  GUIDE
%%   - Restate in your own words; do not paste the problem statement.
%%   - Name every quantity that is asked for, with its symbol and unit.
%%   - Enumerate the questions as Q1, Q2, Q3 (they become the subsections of
%%     Section 6 -- the one-to-one mapping is what judges look for).
%%   - If the statement is ambiguous, state the interpretation you adopt and
%%     why; that sentence is worth more than it looks.
%% =====================================================================
\section{Problem Restatement and Clarification}
\label{sec:problem}

A crewed module of volume $V$ at internal pressure $p_0$ and temperature
$T_0$ develops a circular hole of diameter $d$ in its hull. The external
pressure is $p_a\approx 0$. We are asked to determine:

\begin{enumerate}[label=\textbf{Q\arabic*.},leftmargin=*]
  \item the time $t_{0.3}$ for the internal pressure to fall to
        $0.3\,p_0$, as a function of $d$;
  \item the margin left to a crew of four that needs a given time to
        seal the breach, for $d=\SI{10}{cm}$;
  \item the sensitivity of both times to the parameters that are known
        only to a few per cent.
\end{enumerate}

Two conventions are fixed here. The survival threshold is taken as a total
pressure of $0.3\,p_0$ with an oxygen partial pressure above
\SI{16}{kPa}, following the usual aviation limit; the module is treated as
a single well-mixed volume, so no stratification is resolved. Both choices
are revisited in Section~\ref{sec:discussion}.
